% Expository recovery from the TPK kernel's source owners.
% Provenance and preservation controls are documented in technical/TPK_INTEGRACION.md.
% This file expands the TPK section; it does not replace extension.tex or generacion.tex.

\subsubsection{Dynamic state, observables, and operational dependencies}
\label{tpk-int:estado}

The Topological Phase Kernel organizes the evolution of the two APP sheets
with the orientation determined by TRIT. Its definition comprises the state
on which it acts, the transition rules, and the relations that make
prolongations compatible. Equation \eqref{eq:actualizacion} acquires operational
content when the coordinates received and modified by each factor are specified.
Position belongs to the toroidal support; phase belongs to the calendar;
orientation determines transport; quotients record the arithmetic evaluation
before reduction. These coordinates are related through the transition and
retain different mathematical functions.

The observable projection of a state is
\begin{equation}
 q=(x_+,d_+,x_\times,d_\times,\theta)
 \in\mathcal Q_{\rm obs}
 :=(X_9\times\mathcal D)^2\times\Theta_{108},
 \qquad \mathcal D=\{N,E,S,O\},\quad
 \Theta_{108}=\ZZ/108\ZZ.
 \label{tpk-int:qobs}
\end{equation}
The subscripts distinguish the additive and multiplicative cursors. The
fibers of complete states over \(q\) preserve the ordered trajectory, the
integer evaluations of both sheets, their residues and quotients, and the
orientation and boundary records. A precision prolongation adds the prefix
and the compatible cylinder. Thus two states with the same projection
\(q\) may retain different records.

Operational dependencies have three levels. At the local level, the selector
activates a sheet, transport changes its position, and the update computes
the new arithmetic coordinates. At the trajectory level, these transitions
compose and retain the records of preceding steps. At the refinement level,
extension relations determine which continuation preserves the predecessor's
prefix, signature, and boundary. The finite calendar describes the phase
of these operations; holonomy describes their composition on the history.

\subsubsection{Tritic selection and lifted cardinal transport}
\label{tpk-int:seleccion}

Let \(\bar\theta\in\{0,\ldots,107\}\) be the calendar representative.
The operating phase and the dodecaphase sector are
\begin{equation}
 r(\theta)=1+(\bar\theta\bmod9),\qquad
 \beta(\theta)=\left[\left\lfloor\frac{\bar\theta}{9}\right\rfloor\right]_{12}.
 \label{tpk-int:fase}
\end{equation}
The previously defined selector may be written \(\varepsilon_9=M_{\rm ph}\).
Its value vector in the ordered chart of \(D_9\) is
\[
 m_{\rm ph}=(1,-1,0,1,-1,0,1,-1,0).
\]
The function \(M_{\rm ph}:D_9\to\TRIT\) and this vector are two
presentations of the same selector, with their respective domains.

Define the activation indicators
\[
 a_+(\theta)=\mathbf1_{\{1,4,7\}}(r(\theta)),\qquad
 a_\times(\theta)=\mathbf1_{\{2,5,8\}}(r(\theta)).
\]
Cursor transport is then
\begin{equation}
 x_+'=x_++a_+(\theta)\nu_{d_+},\qquad
 x_\times'=x_\times+a_\times(\theta)\nu_{d_\times}
 \quad\text{in }X_9.
 \label{tpk-int:cursores}
\end{equation}
In phases \(3,6,9\), both positions remain unchanged; the neutral event
forms part of the chronology and its record. Cardinal directions determine
the vectors \(\nu_d\), whereas the phase signature determines which
vector is applied.

The integer lift preserves boundary crossing. If
\(s:X_9\to D_9^2\) chooses positive representatives and
\(x'=x+a\nu_d\), with \(a\in\{0,1\}\), define
\begin{equation}
 b_d(x,a)=\frac{s(x)+a\nu_d-s(x')}{9}\in\ZZ^2.
 \label{tpk-int:borde}
\end{equation}
Divisibility by \(9\) follows from equality in \(X_9\).
If \(z\in\ZZ^2\) records previous crossings, the update is
\(z'=z+b_d(x,a)\). Thus
\[
 s(x')+9z'=s(x)+9z+a\nu_d.
\]
The integer transport data remain reconstructible after the position
has been reduced on the torus.

\begin{proposition}[Conservation of lifted displacement]
\label{tpk-int:prop-desplazamiento}
For a cardinal trajectory of \(n\) steps, with activation indicators
\(a_t\),
\[
 s(x_n)+9z_n=s(x_0)+9z_0+
 \sum_{t=0}^{n-1}a_t\nu_{d_t}.
\]
The same equality is preserved under concatenation of two compatible trajectories.
\end{proposition}
\begin{proof}
The one-step identity is \eqref{tpk-int:borde}. Summing it over the
\(n\) steps cancels the intermediate coordinates. Splitting the summation
interval into two segments produces exactly the concatenation law.
For a closed trajectory, the difference \(z_n-z_0\) is the previously
defined winding number.
\end{proof}

The same principle governs APP evaluations. If a sheet takes the integer
value \(N_t\), its record simultaneously preserves
\(N_t=R_t+9q_t\). The difference between two steps satisfies
\[
 N_{t+1}-N_t=(R_{t+1}-R_t)+9(q_{t+1}-q_t).
\]
Spatial-crossing memory \(z_t\) and arithmetic quotient \(q_t\)
have different definitions; preserving them jointly maintains the
geometric and arithmetic provenance of every transition.

\subsubsection{Two-cursor chronology and composite returns}
\label{tpk-int:cronologia}

The cardinal involution \(\jmath\) interchanges \(N\) with \(S\) and
\(E\) with \(O\). After applying \eqref{tpk-int:cursores}, both
directions are reversed at the end of a block of \(27\) steps.
Writing
\[
 h(\theta)=\mathbf1_{\{0,27,54,81\}}
             ((\bar\theta+1)\bmod108),
\]
the complete rule on the observable projection is
\begin{equation}
 F_{\rm obs}(q)=
 \bigl(x_+',\jmath^{h(\theta)}d_+,
       x_\times',\jmath^{h(\theta)}d_\times,
       \theta+1\bigr).
 \label{tpk-int:fobs}
\end{equation}
The reading that forms the emission uses the positions before transport,
according to the explicit recurrence in Section~\ref{sec:extension}.
The observable update and emission therefore compose in a fixed order.

\begin{proposition}[Reversibility and period of the observable calendar]
\label{tpk-int:periodo}
The map \(F_{\rm obs}\) is a permutation of
\(\mathcal Q_{\rm obs}\) and satisfies \(F_{\rm obs}^{108}=I\).
Within a block of \(54\) steps, positions and directions are restored;
the coordinate in \(\Theta_{108}\) advances by \(54\).
\end{proposition}
\begin{proof}
Given the final state, the preceding phase is first recovered by subtracting
one in \(\Theta_{108}\). That phase determines whether directions were
reversed and which cursor moved. Applying the corresponding involution
and subtracting the cardinal vector uniquely recovers the initial state.

Each interval of \(27\) steps contains nine activations of each cursor.
The following block preserves the activations and reverses the directions,
so the integer displacements cancel over \(54\) steps. The rule has
period \(54\) in its position and orientation increments; the same
cancellation holds for any calendar origin. After \(108\) steps,
\(\theta\) is also restored.
\end{proof}

We reserve
\[
 \mathcal K_{27}=F_{\rm obs}^{27}:
 \mathcal Q_{\rm obs}\longrightarrow\mathcal Q_{\rm obs}
\]
for the iterate of \(27\) transitions. Its fourth iterate is the
identity on the observable state defined in \eqref{tpk-int:qobs}.
Projecting history onto the calendar makes this finite check possible.
The trajectory record, its evaluations, and the new prefixes remain in
the dynamic state on which \(\mathcal U_t\) acts; observable equality
does not reinitialize them.

\subsubsection{Affine memory transport and composition of states}
\label{tpk-int:memoria}

The observable projection admits fibers of arithmetic and geometric data.
Over a configuration \(q\), we write a state in the form
\[
 x=(q,o,h,c,\sigma,C,b,\lambda).
\]
Here \(o\) preserves orientation and sheet; \(h\), residues and
quotients; \(c\), additive transport memory; \(\sigma\), the boundary
signature; \(C\), the precision cylinder; \(b\), its boundary label;
and \(\lambda\), the recorded trajectory. The signature is determined
by a map
\[
 \operatorname{Sig}_q:
 \mathsf H_q\times\mathsf C_q\times\mathsf{Hist}_q
 \longrightarrow\mathsf S_q,
 \qquad \sigma=\operatorname{Sig}_q(h,c,\lambda).
\]
The symbols \(\mathsf H_q,\mathsf C_q,\mathsf{Hist}_q\), and
\(\mathsf S_q\) denote, respectively, the sets of residue--quotient
data, the abelian transport-memory group, the recorded trajectories,
and the boundary signatures. Realizations of this signature are specified
component by component in the regional and terminal constructions.
This dependence prevents treating the signature as a free parameter
once the data producing it are fixed.

The elementary residue--quotient chart is explicit. For
\(n\in\ZZ\), let
\[
 r=1+((n-1)\bmod9),\qquad u=\frac{n-r}{9}.
\]
Then \(n=r+9u\), with \(r\in D_9\) and \(u\in\ZZ\).
Uniqueness follows because two representatives in \(D_9\) congruent
modulo \(9\) are equal. This chart supplies local integer reconstruction.
The precision towers additionally incorporate the truncation maps and
cylinders developed in \ref{sec:extension} and \ref{sec:generacion}.

Let \(\gamma:q\to q'\) be an observable trajectory. Fiber transport
is expressed through maps
\[
 \mathsf H(\gamma):\mathsf H_q\to\mathsf H_{q'},\qquad
 \mathsf C(\gamma):\mathsf C_q\to\mathsf C_{q'},\qquad
 \mathsf{Hist}(\gamma):\mathsf{Hist}_q\to\mathsf{Hist}_{q'}.
\]
The maps respect identities and composition, and
\(\mathsf C(\gamma)\) is an abelian-group homomorphism.
The increment \(k_{\rm tr}(\gamma)\in\mathsf C_{q'}\) satisfies
\begin{equation}
 \begin{split}
 k_{\rm tr}(\operatorname{id}_q)&=0,\\
 k_{\rm tr}(\gamma_2\circ\gamma_1)
 &=k_{\rm tr}(\gamma_2)+
   \mathsf C(\gamma_2)k_{\rm tr}(\gamma_1).
 \end{split}
 \label{tpk-int:cociclo}
\end{equation}
Here \(k_{\rm tr}\) denotes the transport cocycle;
the rational coordinate \(\kappa\) of record \(K\) retains
its notation in the dodecaphase-closure chapter.

Transportable coordinates are updated by
\begin{equation}
 \begin{aligned}
 h'&=\mathsf H(\gamma)h,&
 c'&=\mathsf C(\gamma)c+k_{\rm tr}(\gamma),\\
 \lambda'&=\gamma\circ\lambda,&
 \sigma'&=\operatorname{Sig}_{q'}(h',c',\lambda').
 \end{aligned}
 \label{tpk-int:accion-memoria}
\end{equation}
The second equation is an affine action: the linear term transports
previous memory and the cocycle adds the increment produced by the
trajectory. In the spatial-crossing realization for each cursor separately,
\(\mathsf C(\gamma)=I\) and \(k_{\rm tr}(\gamma)\) is the sum
of the vectors in \eqref{tpk-int:borde}. This concrete instance connects
the abstract law with the cardinal operations already defined.

\begin{proposition}[Compatibility of the update with concatenation]
\label{tpk-int:composicion-memoria}
Updating by \eqref{tpk-int:accion-memoria} along
\(\gamma_1\), followed by updating along \(\gamma_2\),
coincides with updating along \(\gamma_2\circ\gamma_1\).
\end{proposition}
\begin{proof}
For the affine coordinate, two updates give
\[
 c''=\mathsf C(\gamma_2)\mathsf C(\gamma_1)c
       +\mathsf C(\gamma_2)k_{\rm tr}(\gamma_1)
       +k_{\rm tr}(\gamma_2).
\]
Functoriality of \(\mathsf C\) and
\eqref{tpk-int:cociclo} turn this expression into
\(\mathsf C(\gamma_2\circ\gamma_1)c+
k_{\rm tr}(\gamma_2\circ\gamma_1)\).
The statement for \(h\) follows from composition of
\(\mathsf H\), and that for \(\lambda\) from associativity
of trajectories. The final signature agrees because it is evaluated
on the same three reconstructed coordinates.
\end{proof}

The cylinder and its boundary complete this law through the relations
\(\operatorname{Ext}_\gamma\), defined in
\ref{tpk-int:bidireccionalidad}. Their elements determine which of
the preceding updates preserve an admissible prolongation.
States, together with the triples \((\gamma,x,x')\) admitted by
these relations, form a category over observable trajectories:
\[
 (\gamma_2,x',x'')\circ(\gamma_1,x,x')
     =(\gamma_2\circ\gamma_1,x,x'').
\]
The identity of \(x\) is \((\operatorname{id}_q,x,x)\).
Composition preserves arithmetic transport and boundary prolongability
simultaneously. This is the structure in which returns with memory
and successive refinements take place.

\subsubsection{Dodecaphase record, phase, and accumulated memory}
\label{tpk-int:dodecafase}

The \(108\) transitions are distributed over the ordered windows
\[
 E_m=\{9m-8,\ldots,9m\},\qquad 1\leq m\leq12.
\]
The record preserves the phase origin, orientation, and data
produced during each window. Its map is
\[
 \mathcal R_{12}(x)=
 \bigl((E_m,\tau_m,\sigma_m,\kappa_m,\Lambda_m)\bigr)_{m=1}^{12},
\]
where \(\tau_m\) is the window's subroute, \(\sigma_m\) its
orientation, \(\kappa_m\) the integer boundary increment, and
\(\Lambda_m\) the local incidence record. We adopt the chart
developed in \ref{subsec:k-registro-integral}; its window indices
correspond to \(m=\beta+1\).
The domain comprises compatible TPK histories, and the codomain
their oriented temporal records. Composition of segments and their
increments is governed by \eqref{tpk-int:accion-memoria}.
By contrast, the group \(C_{12}\) describes translation of the
window index; \(\Theta_{108}\) preserves the step calendar.
This distinction specifies what is transformed and what is observed
when a cycle is completed.

In coordinates \(\theta=9b+r\), with \(0\leq r<9\), advancing
\(9\) steps preserves \(r\) and increases \(b\) by one.
Addition of phases introduces the quotient
\[
 c_0(r,r')=\left\lfloor\frac{r+r'}9\right\rfloor,
\]
whose cocycle identity ensures associativity of composition.
The finite calendar preserves \([b]_{12}\); a record of turns
preserves \(b\in\ZZ\), or \(b\in\mathbb N_0\) for a
history with fixed origin. Section~\ref{sec:extension} proves
the nonsplit exact sequence \eqref{eq:extension108}. Its function
is to describe the relation between these coordinates before they
are used to construct the signed state.

The signed record receives phase and memory data through the
composition developed in Section~\ref{sec:k-registro}:
\begin{equation}
 R_{\rm sgn}=\Pi_H\circ\mathcal E_{108}^{90,120}
                         \circ\mathcal R_{12}.
 \label{tpk-int:registro-causal}
\end{equation}
Its domain is the terminal state with memory; its output belongs
to \(\ZZ^{12}\). The reversible transformation between this signed
state and \(K\), together with the rational reading of \(K\),
is proved in the corresponding chapter. Record operations precede
these transformations. Reconstructing a record from its cyclic
differences verifies its recoverability; its dynamic provenance
must be preserved through the events that produced it.

For the encoded trace \((d_t)\) retained by the record,
the event reader counts codes \(\{2,4,6,8\}\) in each window
with orientation \((+,+,+,-,-,-,+,+,+,-,-,-)\).
These codes and the cardinal symbols of \(\mathcal D\)
are different record coordinates. Chapter~\ref{sec:k-registro}
develops this reader, the integral decomposition \(1+5+6\),
the congruences characterizing its image, and its exact inverse.
The inversion proof establishes which data this chart preserves;
the definition of \(\mathcal R_{12}\) specifies the oriented
history on which it acts.

\subsubsection{Phase incidence and internal channel construction}
\label{tpk-int:incidencia}

The selector determines three subsets of \(D_9\):
\[
 H_+=\{1,4,7\},\qquad H_-=\{2,5,8\},\qquad H_0=\{3,6,9\}.
\]
Let \(H_{\rm act}=H_+\sqcup H_-\) be the set of active phases
and \(O_{\rm ph}=D_9\setminus\{9\}\) the chart preceding
the marked return. Denote by \(P_2(H_{\rm act})\) the set of
unordered pairs of active phases. Internal incidence is given by
\begin{equation}
 \mathcal F^{\rm ph}_6=H_{\rm act}\times P_2(H_{\rm act}),\qquad
 \mathcal F^{\rm ph}_8=O_{\rm ph}\times P_2(H_{\rm act}),\qquad
 \Theta_{\rm ph}=D_9\times P_2(H_{\rm act}).
 \label{tpk-int:fases90120}
\end{equation}
The indices describe the number of positions in the first component.
The construction uses the phases and calendar origin already fixed;
subsequent exceptional realizations will transport these incidences
to their own domains.

\begin{proposition}[Cardinalities and oriented phase incidence]
\label{tpk-int:cardinales}
The sets in \eqref{tpk-int:fases90120} have cardinalities
\(90,120,135\), respectively. If
\[
 \partial_{\rm ph}=\Theta_{\rm ph}\times\{-1,+1\},\qquad
 E_{\rm ph}=\partial_{\rm ph}\sqcup\{\infty\},\qquad
 L_{\rm ph}=H_0\times P_2(H_{\rm act}),
\]
then
\[
 |\partial_{\rm ph}|=270,\quad |E_{\rm ph}|=271,\quad
 |L_{\rm ph}|=45,\quad
 |\partial_{\rm ph}\sqcup L_{\rm ph}|=315.
\]
Simultaneously translating the phase origin and marked return preserves
all these cardinalities and inclusions.
\end{proposition}
\begin{proof}
There are six active phases, so
\(|P_2(H_{\rm act})|=\binom62=15\). Multiplying by
\(6,8,9\) gives \(90,120,135\). Orientation doubles \(135\);
the return mark adds one element; the three neutral phases give
\(3\cdot15=45\). A calendar translation transports each factor
by a bijection and carries the marked point to its image. Products
and disjoint unions thus preserve the incidence.
\end{proof}

Here the cardinality \(271\) has a phase-incidence realization.
The cubic completion in Section~\ref{sec:extension} has its own
construction of the corona \(1000-729\). The correspondence
between the two realizations must preserve the maps and marks
identifying them, in addition to their cardinality. Likewise,
the selector \(\varepsilon_9\), the incidence \((90,120)\),
and the extractor \(\mathcal E_{108}^{90,120}\) retain the
domains and operations just distinguished.

\subsubsection{Trajectory reversal and conjugate prolongations}
\label{tpk-int:bidireccionalidad}

Bidirectionality begins with an operation on trajectories before scalar
values are assigned to them. For
\(\gamma=(x_0;d_1,\ldots,d_n)\), with endpoint \(x_n\), define
\[
 \operatorname{rev}\gamma
   =(x_n;\jmath d_n,\ldots,\jmath d_1).
\]
The reversed trajectory starts at the endpoint of the original.
Composition satisfies
\[
 \operatorname{rev}^2\gamma=\gamma,\qquad
 \operatorname{rev}(\gamma_2\circ\gamma_1)
   =\operatorname{rev}\gamma_1\circ\operatorname{rev}\gamma_2,
\]
and integer displacement changes sign. These identities follow by
reversing the order of steps and replacing each cardinal vector
with its opposite. For a loop, it follows that
\(\operatorname{wind}(\operatorname{rev}\gamma)
=-\operatorname{wind}(\gamma)\).

On records, reversal also requires transporting the order of components
and the orientation of events. Route reversal and signature conjugation
are distinguishable operations: their composition is specified when
applied to a concrete record. In particular, the regional involution
interchanging closure and propagation components preserves the self-scaling
axis; its formula is given in Section~\ref{mono-ampl:regiones}.
This regional operation acts on a different domain from cardinal
trajectory reversal.

Extensions over a route \(r:q\to q'\) are expressed by
relations between fibers of compatible states. If \(\mathcal F_q\)
denotes the fiber of complete states over \(q\), we require
\[
 \operatorname{Ext}_r\subseteq\mathcal F_q\times\mathcal F_{q'},
 \qquad \Delta_{\mathcal F_q}\subseteq
        \operatorname{Ext}_{\operatorname{id}_q},
\]
\[
 \operatorname{Ext}_{r_2}\circ\operatorname{Ext}_{r_1}
       \subseteq\operatorname{Ext}_{r_2\circ r_1}.
\]
Their composition preserves the evaluations of both sheets, orientation,
quotients, prefix, and boundary. Route reversal does not by itself imply
that every extension relation is a bijection. Conjugate prolongation
is defined on the domain where transport of these data satisfies
the corresponding compatibilities.

\subsubsection{Nonadic holonomy and continuity of refinement}
\label{tpk-int:holonomia}

A phase turn composes transitions on the state with memory:
\begin{equation}
 \Gamma_{9,k}=\mathcal U_{k+8}\circ\cdots\circ\mathcal U_k.
 \label{tpk-int:gamma}
\end{equation}
On an individualized history, this functional composition realizes
the holonomy relation; on the full domain we retain
\(\Gamma_{9,k}\subseteq\widehat X_k\times\widehat X_{k+9}\),
with the admissible prolongations developed later.
The cyclic base has nine phases, while the fiber preserves trajectory,
quotients, and boundary. Monodromy is the action of one turn of this
base on the fibers; holonomy \(\Gamma_{9,k}\) expresses its
realization through the effective TPK transitions.

If \(g\) observes phase and \(M\) the number of retained turns,
\[
 g(\Gamma_{9,k}x)=g(x),\qquad M(\Gamma_{9,k}x)=M(x)+1.
\]
The first identity expresses phase closure and the second identifies
the new temporal depth. Coinductive prolongation also preserves its
cylinders and prefixes. Restrictions between depths compose an inverse
system; a compatible section gathers its antecedents without replacing
them by the latest observation. This arithmetic of histories and its
unity are developed in \ref{hist:seccion}--\ref{hist:doble-varianza};
the regions and their readers are constructed in Section~\ref{sec:generacion}.

Phase transfer on emissions already constructed has a different domain.
Let \(\mathcal A_+\) and \(\mathcal A_\times\) be the
additive and multiplicative signature alphabets, with cardinalities
\(18\) and \(26\). Their product identifies the catalog \(\mathcal U_6\).
For \(f:\mathcal A_+\times\mathcal A_\times\to\RR\),
the sheet averages are
\[
 (E_+f)(s,e)=\frac1{18}\sum_{s'\in\mathcal A_+}f(s',e),\qquad
 (E_\times f)(s,e)=\frac1{26}\sum_{e'\in\mathcal A_\times}f(s,e').
\]
These emissions and their cardinalities are constructed through the
recurrence in Section~\ref{sec:extension}. Once they have been obtained,
the transfer operator is defined on
\(\mathcal U_6\times C_9\) by
\[
 P_+=E_+,\quad P_-=E_\times,\quad P_0=I,\qquad
 \mathcal K_{\rm ph}((u,r),(v,r+1))=P_{M_{\rm ph}(r)}(u,v),
\]
with zero weight for all other phase transitions. The index
\(r\) uses the positive representative of the cyclic class.
The phase order \((+1,-1,0)^3\) then produces the renewal
\[
 \mathcal K_{\rm ph}^{9}\big|_r=E_+E_\times.
\]
Indeed, both averages are idempotent, commute, and average different
factors; their composition assigns weight \(1/(18\cdot26)=1/468\)
to each emission. This equality characterizes a subsequent transfer
observable. History transport remains given by \eqref{tpk-int:gamma},
with its memory and boundary.

The resulting hierarchy establishes the function of the following chapters.
Finite emission constructs the domain and regions; extension relations
preserve their prefixes; the nonadic connection prolongs them; the
dodecaphase record determines the data entering \(K\) and the
closure of \(\alpha\). Numerical and incidence realizations receive
this preceding structure. This order allows every consequence to be
presented in its own chapter while preserving, from the formal kernel
onwards, the objects that make it possible.
